%!TEX root = ../main.tex
\vspace{-8mm}
\section{Introduction}
\label{sec:intro}

% abstract
% Policy-gradient methods are widely used in reinforcement learning, yet training often becomes unstable or slows down as learning progresses. We study this phenomenon through the noise-to-signal ratio (NSR) of a policy-gradient estimator, defined as the estimator variance (noise) normalized by the squared norm of the true gradient (signal). Our main result is that, for (i) finite-horizon linear systems with Gaussian policies and linear state-feedback, and (ii) finite-horizon polynomial systems with Gaussian policies and polynomial feedback, the NSR of the REINFORCE estimator can be characterized exactly—either in closed form or via numerical moment-evaluation algorithms—without approximation. For general nonlinear dynamics and expressive policies (including neural policies), we further derive a general upper bound on the NSR. These characterizations enable a direct examination of how NSR varies across policy parameters and how it evolves along optimization trajectories (e.g., SGD and Adam). Across a range of examples, we find that the NSR landscape is highly non-uniform and typically increases as the policy approaches an optimum; in some regimes it diverges, which can trigger training instability and even policy collapse.

Reinforcement learning (RL) studies how agents learn to make sequential decisions from interaction, and has shown successes across a wide range of domains, including games~\cite{mnih15nature-dqnonatari, silver16nature-alphago, silver17nature-alphazero}, public health~\cite{komorowski18nature-rlmedicine, bastani21nature-rlmedicine}, robotics~\cite{andrychowicz20ijrr-learning, song23sciencerobotics-autonomousracing}, and, more recently, language-model alignment~\cite{ouyang22neurips-rlhf, openai24arxiv-gpt4technicalreport}. At the heart of many advanced RL algorithms stands the \emph{policy-gradient theorem}~\cite{sutton98book-rlbook,sutton99neurips-policygradient,williams92ml-reinforce}, which enables gradient-based learning of policy parameters directly from sampled trajectories, without requiring an explicit model of the environment.

Despite their generality, policy-gradient methods are often brittle: training can become unstable or get stuck~\cite{henderson18aaai-rlimplementation,dohare23ewrl-policycollapse}. 
Prior work has identified multiple reasons for this \emph{instability}. One is the \emph{high variance} of stochastic policy-gradient estimators, which can cause updates to poorly track the true ascent direction; this motivates variance-reduction techniques such as baselines, generalized advantage estimation, and actor--critic methods~\cite{sutton99neurips-policygradient, greensmith04jmlr-variancereduction, schulman15iclr-gae, haarnoja18icml-sac}. 
Another is overly aggressive policy updates, which motivates methods that explicitly constrain update size, \eg via Kullback–Leibler trust regions in TRPO~\cite{schulman15icml-trpo} or clipping in PPO~\cite{schulman17arxiv-ppo}. 
Other related failure modes include \emph{policy collapse} induced by Adam momentum~\cite{dohare23ewrl-policycollapse}, degradation due to \emph{improper baseline selection}~\cite{chung21icml-committalbaselines,mei22neurips-committalbaselines}, and loss of plasticity caused by nonstationary environments~\cite{sokar23icml-dormant, kumar23arxiv-maintaining, muppidi24neurips-trac}.

In this work, we are interested in understanding the mechanisms by which \emph{gradient-estimator statistics} translate into optimization behavior. In particular, we uncover the \emph{non-uniform} structure of policy-gradient variance (\wrt policy parameters) and how it evolves along optimization trajectories.
We focus on the vanilla REINFORCE estimator~\cite{sutton99neurips-policygradient}. Despite its simplicity, it already captures sharp behavior near optima and exposes nontrivial variance effects along optimization trajectories; extensions to richer estimators---including those with baselines or entropy regularization---are discussed in the conclusion.

% \hank{Will this section make the intro too long?}

\textbf{Problem setup and REINFORCE estimator.}
We consider a discounted, finite-horizon Markov decision process with state $s_t \in \Real{n}$, action $a_t \in \Real{m}$, reward $r$, horizon $T$, discount factor $\gamma$, and initial state distribution $\rho_0$.
A trajectory is $\tau := (s_0,a_0,s_1,a_1,\dots,s_{T-1},a_{T-1},s_T)$, generated by a stochastic policy $\pi_\theta(a\mid s)$.
The objective is to maximize the expected return over initial states and trajectories:
$J(\theta) := \mathbb{E}_{s_0\sim\rho_0,\ \tau\sim\pi_\theta}\!\left[R(\tau)\right]$, where
$R(\tau):=\sum_{t=0}^{T-1}\gamma^t\, r(s_{t+1},a_t)$ denotes the discounted return.
We use rewards of the form $r(s_{t+1},a_t)$, instead of $r(s_t,a_t)$, to avoid rewards over a constant initial-state distribution and a terminal action that is sampled but never applied. 
% This help us to consider one-step systems with dynamics envolved.
Following common practice in continuous-control policy optimization, \eg TRPO\,/\,PPO and Linear-Quadratic-Regulator\,/\,Linear-Quadratic-Gaussian~\cite{schulman15icml-trpo,schulman17arxiv-ppo,fazel18icml-lqrglobal},
we focus on Gaussian policies with diagonal covariance:
$
\pi_\theta(a\mid s)=\mathcal{N}\big(\mu_\theta(s),\,\Sigma\big),
\Sigma=\mathrm{diag}(\sigma^2),
\ell=\log\sigma\in\mathbb{R}^m.
$
We restrict attention to the finite-horizon setting because (i) many widely used benchmarks are episodic~\cite{brockman16arxiv-gym}, and (ii) stability considerations for structured systems (\eg linear--quadratic control under linear Gaussian policies) can be handled cleanly. We consider the REINFORCE policy gradient estimator
\begin{align}\label{eq:policy_gradient_theorem_1}
\nabla_\theta J(\theta)\!=\!\mathbb{E}\!\left[\widehat G_\theta\!:=\!\left(\sum_{t=0}^{T-1} \nabla_\theta \log \pi_\theta(a_t\mid s_t)\right)\!R(\tau)\right]\!\!.\!\!
\end{align}

\vspace{-2mm}
\textbf{Noise-to-signal ratio (NSR).}
To relate gradient-estimator noise to optimization behavior, in addition to the variance of the gradient estimator, we study the \emph{noise-to-signal ratio} (NSR),
defined as the estimator variance (noise) normalized by the squared norm of the true gradient (signal):
% \footnote{We use Frobenius norm throughout to unify notation.}:
\[
\mathrm{NSR}(\theta)
:=\frac{\mathrm{Var}_{\mathrm{Fro}}\big(\widehat G_\theta\big)}{\|\nabla J(\theta)\|_{\mathrm{F}}^2},
\ \ 
\mathrm{Var}_{\mathrm{Fro}}(\widehat G)
:=\mathbb{E}\big[\|\widehat G-\mathbb{E}[\widehat G]\|_{\mathrm{F}}^2\big].
\]
%\cite{solodov98compopt-stronggrowth,cevher19optletters-sgdlinearconvergence,mccandlish18arxiv-batchsize,bottou18siamreview-sgd,byrd12mp-batchsize}.
To avoid degeneracy at stationary points, we report NSR only for $\|\nabla J(\theta)\|_{\mathrm{F}}>0$. NSR is a natural measure of how informative a stochastic gradient is relative to its mean, and is closely related to the \emph{strong growth condition} underlying favorable convergence rates and batch-size rules in stochastic optimization~\cite{bottou18siamreview-sgd, cevher19optletters-sgdlinearconvergence}.
% \cyan{While NSR can be estimated by Monte Carlo rollouts, doing so can be expensive and offer limited insight into \emph{where} in parameter space the estimator becomes uninformative.} \red{--- putting this sentence here is awkward, see where I added it.} 

% \textbf{Contributions.}
% Classical analyses of stochastic optimization often assume uniformly bounded gradient variance or NSR.
% Our central contribution is to show that even in structured systems, the NSR of the REINFORCE estimator can be highly non-uniform across parameter space and can increase sharply as the policy approaches an optimum.

% On the theory side, we develop exact characterizations of the REINFORCE mean and variance (hence NSR) under Gaussian policies in linear--quadratic control, including (i) closed-form expressions in the one-step case that reveal explicit dependence on the dynamics, feedback gain, initial-state covariance, and policy covariance; and (ii) a finite-horizon treatment via lifted dynamics and Gaussian moment evaluation that yields an exact numerical procedure for computing NSR.
% We further extend the same moment-evaluation machinery to polynomial dynamics with polynomial feedback, and we derive variance upper bounds that apply to general nonlinear dynamics and neural policies.
% Through these characterizations, we show that NSR scales inversely with the policy covariance as the policy becomes more deterministic, and scales proportionally with the initial-state covariance as the initial distribution becomes broader.
% Moreover, for linear systems we prove that NSR can grow exponentially with the horizon when the closed-loop dynamics are unstable.

% Experiments on both systems show that along optimization trajectories (GD/SGD/Adam), NSR typically \textbf{increases} as the policy approaches an optimum, and gradient descent (GD) can reach low-covariance (nearly deterministic) solutions, whereas SGD often fails to converge as NSR grows.
% To interpret this behavior, Corollary~\ref{cor:lqg_stationary} implies that the optimal LQG policy is deterministic (i.e., $\Sigma\to \mathbf{0}$ at stationarity).
% Consequently, as the policy approaches optimality, the NSR can remain large or even diverge, providing a concrete mechanism for slowdowns and instability near optima.
% This highlights an inherent exploration--exploitation tension for REINFORCE-style updates.

% We summarize our main contributions as follows:
% \vspace{-3mm}
% \begin{itemize}
%     \item \textbf{Exact/computable NSR.} Exact NSR formulas for linear systems, extensions to polynomial systems, and variance upper bounds for nonlinear/neural policies.
%     \item \textbf{Scaling laws.} NSR increases as $\Sigma$ decreases and as $\Sigma_0$ increases, and can grow exponentially with horizon under unstable closed-loop dynamics.
%     \item \textbf{Optimization implication.} Since optimal LQG policies are deterministic (Corollary~\ref{cor:lqg_stationary}), approaching optimality can drive NSR upward, explaining slowdowns/instability and why GD can reach near-deterministic solutions while SGD often struggles.
% \end{itemize}
% \vspace{-3mm}

\textbf{Contributions.}
Classical analyses of stochastic optimization often assume \emph{uniformly bounded} gradient variance or NSR.
Our central contribution is to show that, in structured RL problems, the NSR and the variance of the REINFORCE estimator can be highly non-uniform and can increase sharply---potentially grow unbounded---as the policy approaches optimality. Our theoretical results include:

\vspace{-3mm}

\begin{itemize}
    \item (Theorem~\ref{thm:lqg-variance-one-step}) \textbf{Closed-form expressions of the NSR for one-step linear systems}, 
    revealing explicit dependence of the NSR on the dynamics, feedback gain, initial-state covariance, and policy covariance.
    \item (Theorems~\ref{thm:lqr-variance-T-compose} \&~\ref{thm:lqr-variance-T-bound}) \textbf{Exact computation of the NSR for multi-step linear systems.} We provide an exact (non-asymptotic) procedure to compute the NSR in finite-horizon linear systems with linear-Gaussian policies via lifted dynamics and Gaussian moment evaluation, together with bounds that expose key scaling factors.
    \item (Proposition~\ref{prop:poly-isserlis}) \textbf{Exact computation of the NSR for polynomial systems.}
    This is an extension of the moment-evaluation approach to polynomial dynamics with polynomial feedback under Gaussian exploration.
    \item (Theorem~\ref{thm:nonlinear-bound}) \textbf{Upper bounds for nonlinear systems.}
    For general nonlinear dynamics and expressive (including neural) Gaussian policies, we provide an upper bound on the variance of the REINFORCE estimator.
\end{itemize}

\vspace{-4mm} 
Through these characterizations, we show that the variance and the NSR of the REINFORCE estimator scale inversely with policy covariance as the policy becomes more deterministic, and scale proportionally with the initial-state covariance as the initial distribution becomes broader. Moreover, we prove they can grow exponentially with the horizon when the closed-loop dynamics are unstable for linear systems.

These theoretical results enable us to numerically investigate how the NSR evolves along optimization trajectories. We show---for both linear and polynomial systems---that the NSR typically \emph{increases} as the policy approaches optimality. Moreover, gradient descent (GD) can drive the policy toward nearly deterministic solutions, whereas stochastic gradient descent (SGD) often fails to converge as the NSR grows. This behavior is consistent with Corollary~\ref{cor:lqg_stationary}, which implies that the optimal policy is deterministic (\ie the policy covariance satisfies $\Sigma\to \mathbf{0}$). Consequently, as the policy nears optimality, the NSR can grow unbounded, providing a concrete mechanism for slowdowns and instability near optima. This highlights an inherent exploration--exploitation tension for REINFORCE-style updates.

\vspace{-1mm}
We stress the value of \emph{exact} characterizations of REINFORCE variance and the NSR. Although both can be estimated from Monte Carlo rollouts, accurate estimation often requires a prohibitively large number of trajectories; with only a modest rollout budget, sampling noise can yield variance and NSR estimates that misrepresent policy-gradient optimization dynamics, as we illustrate in Appendix~\ref{sec:app-add-experiments}.

\vspace{-1mm}
\textbf{Paper organization.}
In Section~\ref{sec:lqg}, we characterize the NSR for linear--quadratic systems under Gaussian policies; we then extend the analysis to polynomial systems (Section~\ref{sec:poly}) and nonlinear dynamics (Section~\ref{sec:nonlinear_new}). Each section presents theory first, then numerical experiments.
All the proofs and related work not covered above are deferred to the appendix.




% Reinforcement learning (RL)~\cite{schulman17arxiv-ppo,schulman15icml-trpo} has achieved notable successes in robotics and is increasingly used in natural language processing~\cite{ouyang22neurips-rlhf}. \hy{This is not a robotics paper, we need to downplay robotics. The major success of RL before robotics is in games like AlphaGo and even medical applications like AlphaFold? We should be more general here about the applications of RL.}
% At the heart of many advanced RL algorithms stands the \emph{policy-gradient lemma}~\cite{sutton98book-rlbook,sutton99neurips-policygradient}, \hy{which allows gradient-based learning of optimal policies purely from environment interactions without explicit model of the environment.} However, policy-gradient training can be brittle due to high estimator variance. Despite standard variance-reduction techniques such as baselines, GAE \hy{spell out GAE}, and actor--critic methods~\cite{greensmith04jmlr-variancereduction, schulman15iclr-gae, sutton99neurips-actorcritic, haarnoja18icml-sac}, the mechanisms behind this brittleness are not fully understood. In particular, comparatively little work connects the optimization behavior of policy-gradient methods to the statistical structure of their gradient estimators.

% \textbf{Problem setup.} We consider a finite-horizon Markov decision process with state $s \in \Real{n}$, action $a \in \Real{m}$, and trajectories
% $\tau := (s_0,a_0,s_1,a_1,\dots,s_T,a_T)$ \hy{Do we need $a_T$?}
% generated by a stochastic policy $\pi_\theta(a\mid s)$ parametrized by $\theta\in\mathbb{R}^d$. The learning objective is to maximize the expected return $R(\tau)$ \hy{over a distribution of the initial state $s_0$}. Following common practices in modern RL---including TRPO~\cite{schulman15icml-trpo}, PPO~\cite{schulman17arxiv-ppo}, and LQR/LQG policy optimization~\cite{fazel18icml-lqrglobal}---we focus on Gaussian policies for continuous control:
% $
% \pi_\theta(a\mid s)=\mathcal{N}\big(\mu_\theta(s),\,\Sigma\big),~\Sigma=\mathrm{diag}(\sigma^2),~\ell=\log\sigma\in\mathbb{R}^m.
% $
% We restrict our attention to the finite-horizon setting because it matches the widely used benchmark environments such as Gym~\cite{brockman16arxiv-gym}, and (ii) it simplifies stability considerations for linear--quadratic systems under linear gaussian policies.


% \hy{
%     The structure of the introduction is a bit confusing. There is also redundancy with the preliminaries. I suggest we merge the preliminaries into the introduction and do the following
%     \begin{itemize}
%         \item Start with the background that RL is very general and successful, and at the heart of modern RL algorithms lies the policy gradient lemma, which enables gradient-based learning purely from trajectories.
%         \item A paragraph on ``Instability'': however, RL training using policy gradients is known to be unstable and brittle (give some references and examples). There has been several lines of works studying the reasons behind this instability --> summarize related works. (i) TRPO and PPO argue that instability comes from overly large policy updates and therefore encourage small updates using KL divergence or clipping. (ii), (iii) etc. 
%         \item Start a new paragraph: however, comparatively little work connects the optimization behavior of policy-gradient methods to the statistical
%         structure of their gradient estimators. In this work, our main focus is to investigate the non-uniform structure of the variance of the gradient estimator. Towards this end, we start by presenting the problem setup and the REINFORCE estimator. Say that we focus on the REINFORCE estimator because the structure of the REINFORCE estimator is already rich enough, and we leave the extension to other estimators, particularly with baselines, for future work.
%         \item A paragraph on ``Problem setup and REINFORCE estimator'': the usual MDP setup, the policy setup, the finite-horizon setting, finally give the REINFORCE estimator.
%         \item A paragraph on ``Noise-to-signal ratio (NSR)'': define NSR, and explain why we want to study it: first, the NSR is a natural quantity that measures how noisy the gradient estimator is compared with the true gradient itself; second, the NSR is closely related to the strong growth condition, under which one can prove favorable convergence rates for stochastic optimization.
%         \item A paragraph on ``Contributions'': while classical literature in stochastic optimization assumes the variance of the gradient estimator is uniformly bounded or the NSR is uniformly bounded, the central contribution of this work is to show the non-uniform structure of the NSR and how that relates to the optimization behavior of policy-gradient methods...
%     \end{itemize}
% }

% \textbf{REINFORCE estimator.} A standard Monte Carlo estimator for the policy gradient using $N$ sampled trajectories yields the REINFORCE estimator~\cite{sutton99neurips-policygradient}. Its practical performance is often governed by estimator variance, which can be large and can lead to unreliable stochastic gradient \emph{descent} (SGD) updates. TRPO argues that overly large policy updates can ruin performance and therefore imposes a Kullback--Leibler (KL) divergence constraint to limit step size. Dohare et al.~\cite{dohare23ewrl-policycollapse} study \emph{policy collapse}, where performance suddenly degrades during training. Closely related to our variance study, Mei et al.~\cite{mei22neurips-committalbaselines} show that the variance of policy-gradient estimators can become \emph{unbounded} near optima in certain bandit settings.

% \textbf{Noise-to-signal ratio (NSR).} To connect gradient-estimator noise to SGD behavior, we study the \emph{noise-to-signal ratio} (NSR), defined as the variance of a stochastic gradient estimator divided by the squared norm of the true gradient. NSR is closely related to the strong growth condition~\cite{solodov98compopt-stronggrowth, cevher19optletters-sgdlinearconvergence, mccandlish18arxiv-batchsize}, and it appears in convergence guarantees and batch-size selection rules for SGD~\cite{bottou18siamreview-sgd, schmidt13arxiv-sgdconveregnce, byrd12mp-batchsize}. Although NSR can be estimated by Monte Carlo rollouts, doing so can be expensive and offers limited insight into how NSR scales with problem parameters.

% \textbf{Contributions.} We first derive exact expressions for the NSR of the REINFORCE estimator in one-step linear--quadratic systems with Gaussian policies. The resulting formulas depend explicitly on the system dynamics matrices, the linear feedback gain, the initial-state covariance, and the policy covariance. We show that NSR scales inversely with policy variance as the policy becomes more deterministic, and scales proportionally with the initial-state variance as the initial distribution becomes broader. We then extend our analysis to finite-horizon linear--quadratic systems with Gaussian policies. While closed-form expressions are generally unavailable, we provide an exact procedure to compute NSR via recursive Gaussian moment evaluation. We also derive an upper bound on the estimator variance and show that NSR retains the same scaling with respect to initial-state covariance and policy variance. Moreover, we prove that NSR can grow exponentially with the horizon when the closed-loop dynamics are unstable. \red{The one-step analysis also yields insights relevant to actor--critic methods.}

% We further extend the moment-evaluation approach to polynomial dynamics with polynomial policies, where we again obtain an exact NSR characterization. Experiments on both linear and polynomial systems show that along optimization trajectories (GD/SGD/Adam), NSR typically increases as the policy approaches an optimum; correspondingly, SGD and Adam tend to slow down as the policy becomes more deterministic. Corollary~\ref{cor:lqg_stationary} implies that the optimal policy is deterministic (i.e., the policy covariance satisfies $\Sigma\to \mathbf{0}$). Consequently, as the policy approaches optimality, the NSR can remain unbounded, providing a concrete mechanism for slowdowns and instability near optima. This highlights an inherent exploration--exploitation tension for REINFORCE-style updates in LQG. Our experiments corroborate the theory: gradient descent (GD) can reach low-covariance (nearly deterministic) solutions, whereas SGD often fails to converge due to an increasingly large NSR. 


% By Monte Carlo sampling, the expected return can be approximated as the mean of $N$ sampled trajectories. In practice, practitioners have long found policy-gradient methods to be fragile and sensitive on difficult tasks. Trust Region Policy Optimization (TRPO) \cite{schulman15icml-trpo} argues that large steps can collapse performance and imposes a Kullback–Leibler (KL) divergence constraint to limit step size. Dohare et al. \cite{dohare23ewrl-policycollapse} study the phenomenon of policy collapse, where performance suddenly degrades during training. They attribute this to abrupt changes in Adam’s momentum, which in turn lead to a ``deadly trial'' effect with off-policy data. Moalla et al. \cite{moalla24neurips-rlrepresentation} investigate the impact of representation collapse on policy-gradient methods, showing that rank-deficient representations can yield brittle policies. More closely, Chung et al. \cite{chung21icml-committalbaselines} show that even when the variance is held fixed, different choices of baseline can lead to markedly different convergence behavior. They introduce the notion of committal behavior to characterize when a policy tends to ``doubt itself.'' Building on this, Mei et al. \cite{mei22neurips-committalbaselines} argue that, rather than focusing solely on variance reduction, one should also account for the aggressiveness of the parameter updates induced by different baselines.
